\subsubsection{Source identity, review and storage}

The source store preserves original bytes and an exact text derivative. The raw
document and derivative each have a SHA-256 digest. A source span records the
derivative hash, page, half-open character offsets and a hash of the quoted text.
Extraction changes create a new derivative; they do not silently move an approved
span. A reviewer can open the original page beside the extracted passage.
Scanned or poorly extracted clauses require visual confirmation before release.

The source register distinguishes circulars, annexes, FAQs, codes, bank policy
and interpretation. Dependency rows record references, definitions, amendments
and replacements; unresolved targets remain explicit. Publication date does not
automatically establish effectiveness. Applicability records identify legal
entity, regulated role, activity, product, client class and jurisdiction. A
product-provider requirement cannot enter the bank's distributor rules merely
because its title mentions a product the bank sells.

The bundle contains definitions, expressions, spans and interpretations. Its
canonical JSON hash identifies the entire immutable revision. Review decisions
live outside those bytes and approve that exact hash. The lifecycle is draft,
in review, approved, released and retired; rejection returns a revision to draft,
while an edit creates a new draft hash. Release requires distinct compliance and
engineering reviewers, resolved blocking questions and dependencies, a supported
effective interval and passing evidence for that hash. The author cannot approve
their own revision under the proposed prototype policy.

SQLite transactions enforce review revisions and idempotency. A stale review
update is rejected rather than overwriting another review. Content-addressed
blobs are written before their database references; a missing referenced blob
is an integrity error. Audit rows record accepted transitions. Hashes detect
accidental corruption but do not make a database tamper-proof against an
administrator able to rewrite both records and hashes. Enterprise identity,
signed retention and access controls belong to the later bank deployment.

\subsubsection{The small financial fixture within the complete dry run}

Section~\ref{sec:dry-run} follows the complete selected transaction profile.
The smaller conformance fixture imports 23EC35 and Annex 1, verifies their hashes and
opens paragraph 3.1 beside the financial rule. Its two wealth definitions are
review issues until the relevant source and data mappings are accepted. An
engineer loads the supplied draft bundle and the synthetic \texttt{f01} case:
portfolio HK\$40 million and qualifying net assets zero. The trace evaluates
both inclusive comparisons and their alternative, returning \texttt{TRUE} for
\texttt{spi.financial}. It also identifies the exact source revision, input
snapshot and engine version.

The reviewer next opens \texttt{f02}, one cent below the portfolio threshold,
with the other route still false. This returns \texttt{FALSE}. In \texttt{f07},
the portfolio is unknown and net assets are HK\$70 million, so the answer is
\texttt{UNKNOWN}. In \texttt{f06}, changing only net assets to HK\$90 million
establishes the financial condition despite the still-missing portfolio. These
cases let a compliance reader inspect the interpretation before any system rule
is exported. A proposed strict comparison is challenged with \texttt{f01}.

Once dependencies and meanings are reviewed, the evidence report is attached to
the bundle hash, reviewers approve it and the release service exports the package.
The production change package includes deterministic Java source and its tested
JAR against the same decision contract. A host adapter supplies the bank's
facts and enforces the released-version and concurrency requirements.
A changed rule creates another hash and invalidates release approval. Historic
decisions still point to the old hash and their original evidence.

\subsubsection{Service boundaries an implementer can use}

The API imports sources, creates and retrieves bundles, records reviews, releases
approved hashes, evaluates snapshots, appends events, replays histories and
compares versions. The transport contract is summarized below; companion files
preserve the same field names for implementation. The central reference call is

\begin{minipage}{\textwidth}
\begin{lstlisting}
evaluate(bundle, snapshot, rule_id, valid_at, known_at,
         mode="draft")
    -> EvaluationResult
\end{lstlisting}
\end{minipage}

Its response includes the tagged outcome, exact value if any, reason codes,
missing and blocking inputs, trace, source/bundle/input hashes and engine version.
Production calls require a released exact hash; draft calls visibly retain draft
mode. An HTTP failure is different from a well-formed request whose assessment
returns an error or conflict.

The local authoring service uses the following endpoints. Path parameters name
immutable content identities; state-changing requests also carry a caller-scoped
idempotency key. Repeating the same key and request returns the stored response;
reusing it with different content returns HTTP 409. Optimistic revisions prevent
two reviewers from overwriting each other's changes.
\begingroup\small
\begin{longtable}{P{0.36\textwidth}P{0.54\textwidth}}
\caption{Authoring-service contract. This service is separate from the Java
transaction library.}\\
\toprule Method and path & Required request and returned result \\\midrule
\endfirsthead
\toprule Method and path & Required request and returned result \\\midrule
\endhead
\texttt{POST /v1/sources/import} & Official URL or local blob, reference ID and media type; returns source revision and derivative hashes. Invalid content is 422; fetch failure is 502.\\
\texttt{POST /v1/bundles} & Complete typed bundle; returns immutable hash, draft status and revision 1. Schema or type failure is 422.\\
\texttt{GET /v1/bundles/\{hash\}} & Returns exact bundle, review state and issues; unknown hash is 404.\\
\texttt{POST /v1/bundles/\{hash\}/review} & Decision, reviewer identity/role, comment and expected revision; returns new state/revision. Stale update is 409; invalid transition is 422.\\
\texttt{POST /v1/bundles/\{hash\}/release} & Evidence-report hash, Java-release-manifest hash and expected revision; returns release record if all guards pass. Missing approval or mismatched code evidence is 422.\\
\texttt{POST /v1/evaluations} & Bundle hash, snapshot, rule ID, mode and both times; returns the complete evaluation. A nonreleased production bundle is 409.\\
\texttt{GET /v1/evaluations/\{hash\}} & Returns preserved request, result and source links; unknown hash is 404.\\
\texttt{POST /v1/event-streams/\{id\}/events} & Event, expected revision and ordering authority; returns stored event/revision. An ID payload mismatch or sequence collision is 409.\\
\texttt{POST /v1/event-streams/\{id\}/replay} & Both times and profile version; returns immutable consent/obligation state with evidence.\\
\texttt{POST /v1/comparisons} & Old/new bundle hashes, rule ID and declared domain; returns a persisted job ID and later a counterexample or a specifically bounded analysis result.\\
\texttt{GET /v1/changes/\{old\}/\{new\}} & Returns changed definitions, rules, sources, affected dependants and evidence that must be regenerated.\\
\bottomrule
\end{longtable}
\endgroup

The persistent relationships follow the questions the reviewer asks. A source
revision points to immutable raw and extracted blobs; an interpretation points
to its spans and unresolved issues; a bundle contains the interpreted rules;
reviews and release records point to that bundle hash. A snapshot identifies
its subject and facts. An evaluation points to the bundle, snapshot and both
times. An event belongs to a stream and has an immutable identifier and ordering
evidence; replay results are separate records. Jobs, idempotency records and audit
events preserve asynchronous work and accepted state transitions. The supplied
SQL represents these relationships with foreign keys and JSON payloads; a Java
release manifest can be retained as an immutable blob referenced from the
release evidence. No mutable row is allowed to change the contents identified
by an old digest.

For content identity, sort JSON object keys by ASCII order, preserve array order
and string code points, use compact UTF-8 and hash the bytes with SHA-256.
Reject duplicate keys, floating-point values, nonfinite values and unpaired
Unicode surrogates. JSON integer metadata is limited to the exactly portable
range $[-(2^{53}-1),2^{53}-1]$; domain integers and money use decimal strings.
This is an explicitly specified local format, not an assertion of conformance
to a different canonical-JSON standard. A result hash excludes itself and
incidental timing/request identifiers. These rules let the Python authoring
service and Java library identify the same inputs and semantic output.

Longer solver or drafting calls use persisted jobs with explicit failed,
cancelled and timed-out states. They cannot alter approval records. Official
documents and model output remain untrusted input: fetching restricts hosts and
redirects, validates content and limits resources; parsing never executes a
document's instructions. The initial service binds to localhost and uses public
sources and synthetic facts. This makes the proposed behavior implementable
before private data access or production integration is considered.
